\chapter{Modellvergleich und Reflexion}
\label{chap:modellvergleich}

\begin{myremark}[Algorithmusart in dieser DL]
Funktionsart: \textbf{Meta-Vergleich} über Klassifikation, Regression, Assoziation und Filterung. \\
Umsetzungsart: \textbf{vergleichend} zwischen regelbasiert und datengetrieben.
\end{myremark}

\begin{myremark}
Diese DL fasst die behandelten Modelle zusammen und schärft den Blick für
die \textbf{richtige Modellwahl} sowie gesellschaftliche Aspekte.
\end{myremark}

\section{Vergleich der Methoden}

\begin{table}[H]
\centering
\small
\begin{tabular}{p{2.8cm}p{2.1cm}p{3.5cm}p{3.5cm}}
\toprule
\textbf{Modell} & \textbf{Aufgabe} & \textbf{Stärken} & \textbf{Grenzen} \\
\midrule
Entscheidungsbaum & Klassifikation & sehr gut interpretierbar; keine Normalisierung nötig & Overfitting-anfällig ohne Beschränkung \\
\midrule
Random Forest & Klassifikation & robust, oft gute Genauigkeit, Feature Importance & weniger gut erklärbar; langsamer \\
\midrule
SVM & Klassifikation & effektiv bei klar trennbaren Klassen, Kernel-Trick & Parameterwahl anspruchsvoll; langsam bei grossen Daten \\
\midrule
KNN & Klassifikation / Regression & einfach, intuitiv, kein Training & sensitiv auf Skalierung; langsam bei Vorhersage \\
\midrule
Lineare Regression & Regression & schnell, gut interpretierbar & nur lineare Zusammenhänge \\
\midrule
Neuronales Netz & Beides & sehr leistungsfähig bei komplexen Daten & viel Daten nötig; Black Box; hoher Rechenaufwand \\
\midrule
Apriori & Assoziation & findet gut verständliche Muster in Transaktionen & viele triviale Regeln ohne Schwellwerte \\
\midrule
Naive Bayes & Filterung / Klassifikation & schnell, robust bei Textklassifikation & Unabhängigkeitsannahme oft verletzt \\
\bottomrule
\end{tabular}
\caption{Überblick über die behandelten Verfahren}
\label{tab:modellvergleich-15}
\end{table}

\section{Das richtige Modell wählen}

\begin{myexercise}[Modellwahl begründen]
Welches Modell würden Sie in folgenden Situationen einsetzen?
Begründen Sie jeweils kurz.
\begin{enumerate}
    \item Ein Arzt möchte nachvollziehen, warum das Modell eine Diagnose stellt.
    \item Ein Social-Media-Unternehmen klassifiziert täglich Millionen von Posts.
    \item Sie haben 200 sauber gelabelte Datenpunkte und wenig Rechenzeit.
    \item Sie möchten den Stromverbrauch nächster Woche vorhersagen.
\end{enumerate}
\begin{myanswer}
1.\ Entscheidungsbaum (hohe Interpretierbarkeit). \\
2.\ Random Forest oder SVM (skalierbar, robust). \\
3.\ KNN oder Entscheidungsbaum (einfach, schnell trainiert). \\
4.\ Lineare Regression (Regression, gut interpretierbar).
\end{myanswer}
\end{myexercise}

\section{Overfitting und Underfitting – Wiederholung}

\begin{myexercise}[Lernkurven interpretieren]
Ein Modell zeigt folgende Genauigkeiten in Abhängigkeit der Trainingsgrösse:
\begin{itemize}
    \item Bei 50 Trainingspunkten: Train\,=\,98\,\%, Test\,=\,61\,\%
    \item Bei 500 Trainingspunkten: Train\,=\,89\,\%, Test\,=\,85\,\%
\end{itemize}
\begin{enumerate}
    \item Welches Modell zeigt Overfitting, welches ist besser generalisiert?
    \item Was empfehlen Sie: mehr Daten oder ein einfacheres Modell?
\end{enumerate}
\end{myexercise}

\section{Gesellschaftliche Aspekte}

\begin{myexercise}[Bias in ML-Modellen]
\begin{enumerate}
    \item Erklären Sie, wie \textbf{Datenbias} entsteht und nennen Sie ein Beispiel.
    \item Ein Kreditmodell wurde vorwiegend auf Daten von Männern trainiert.
          Welche Konsequenzen kann das für Frauen haben?
    \item Welche Massnahmen könnten Bias verringern?
\end{enumerate}
\end{myexercise}

\section{Lernziele des Skripts}

\begin{todolist}
    \item Ich kann den Unterschied zwischen \textbf{Klassifikation} und \textbf{Regression} erklären.
    \item Ich kann die Grundidee von \textbf{Entscheidungsbaum}, \textbf{Random Forest},
            \textbf{SVM}, \textbf{KNN}, \textbf{Regression}, \textbf{Apriori},
            \textbf{Naive Bayes} und \textbf{Neuronalen Netzen} beschreiben.
    \item Ich kann \textbf{Overfitting} erkennen und geeignete Massnahmen nennen.
    \item Ich kann ein Modell in Python trainieren und eine Vorhersage mit geeigneten
          Metriken einordnen.
    \item Ich kann für ein gegebenes Problem das passende Modell auswählen und begründen.
    \item Ich kann Bias-Risiken benennen und Gegenmassnahmen vorschlagen.
\end{todolist}
